\documentclass[11pt]{article}
\usepackage[margin=1in]{geometry}
\usepackage{xcolor}
\usepackage[breakable,listings]{tcolorbox}
\tcbuselibrary{listings}
\usepackage{hyperref}
\usepackage{enumitem}
\usepackage{listings}
\usepackage{verbatim}
\usepackage{amsmath}
\usepackage{graphicx}
\usepackage{booktabs}
\usepackage{float}

% Define colors
\definecolor{osaorange}{RGB}{220,115,38}
\definecolor{aiblue}{RGB}{30,144,255}
\definecolor{codebackground}{RGB}{248,248,248}

% Configure hyperref
\hypersetup{
    colorlinks=true,
    linkcolor=blue,
    urlcolor=blue,
    citecolor=blue
}

% Configure listings for code display
\lstset{
  basicstyle=\small\ttfamily,
  breaklines=true,
  breakatwhitespace=true,
  postbreak=\mbox{\textcolor{gray}{$\hookrightarrow$}\space},
  columns=flexible,
  keepspaces=true,
  showstringspaces=false
}

% Code environment using tcolorbox with listings
\newtcblisting{rcode}{
  colback=white,
  colframe=gray!50,
  boxrule=0.5pt,
  arc=0pt,
  left=3pt,
  right=3pt,
  top=3pt,
  bottom=3pt,
  width=\textwidth,
  breakable,
  listing only,
  listing options={
    basicstyle=\small\ttfamily,
    breaklines=true,
    breakatwhitespace=true,
    postbreak=\mbox{\textcolor{gray}{$\hookrightarrow$}\space},
    columns=flexible,
    keepspaces=true,
    showstringspaces=false
  }
}

\title{}
\author{}
\date{}

\begin{document}

% Header box
\begin{tcolorbox}[colback=osaorange,colframe=black,boxrule=2pt,arc=0pt,left=3pt,right=3pt,top=6pt,bottom=6pt,boxsep=2pt]
\centering
\textcolor{white}{\textbf{\Large H524 Introduction to Biostatistics}}\\
\textcolor{white}{\textbf{\Large Week 6 Lab: Power, Two-Sample Tests, and ANOVA}}\\
\textcolor{white}{\textbf{\Large Fall 2025}}
\end{tcolorbox}

\vspace{8pt}

\noindent\textbf{Format:} Online\\
\textbf{Tools Required:} R, RStudio\\
\textbf{Estimated Time:} 90 minutes

\section{Learning Objectives}

By the end of this lab, you will be able to:
\begin{enumerate}
\item Calculate statistical power and required sample sizes in R
\item Perform independent two-sample t-tests (pooled and Welch's)
\item Conduct paired t-tests for dependent data
\item Run one-way ANOVA to compare multiple groups
\item Perform post-hoc tests (Tukey's HSD) after significant ANOVA
\item Check assumptions for t-tests and ANOVA
\item Interpret output from hypothesis tests correctly
\item \textcolor{aiblue}{\textbf{Use AI tools to help with power analysis and test selection}}
\end{enumerate}

\section{Part 1: Statistical Power and Sample Size}

Statistical power is the probability of correctly rejecting the null hypothesis when it is false. Before conducting a study, we should calculate the required sample size to achieve adequate power (typically 0.80 or 80\%).

\subsection{Exercise 1.1: Understanding Power}

\textbf{Scenario:} We want to test if a new drug reduces cholesterol. Population mean is 200 mg/dL with SD = 30 mg/dL. We want to detect a reduction to 190 mg/dL.

\begin{rcode}
# Power calculation for one-sample t-test
# Parameters
mu0 <- 200          # Null hypothesis mean
mu1 <- 190          # True mean (alternative)
sigma <- 30         # Population SD
n <- 25             # Sample size
alpha <- 0.05       # Significance level

# Manual calculation
se <- sigma / sqrt(n)  # Standard error = 6

# Critical values for two-sided test
z_crit <- qnorm(1 - alpha/2)  # 1.96
lower_crit <- mu0 - z_crit * se  # 188.24
upper_crit <- mu0 + z_crit * se  # 211.76

cat("Rejection region:", lower_crit, "to", upper_crit, "\n")

# Power: P(reject H0 | mu = 190)
# We reject if sample mean falls in rejection region
power_lower <- pnorm(lower_crit, mean=mu1, sd=se)
power_upper <- pnorm(upper_crit, mean=mu1, sd=se,
                     lower.tail=FALSE)
power <- power_lower + power_upper

cat("Power:", round(power, 3), "\n")
cat("Type II error (beta):", round(1-power, 3), "\n")

# Interpretation: Only 34.5% chance of detecting the effect!
\end{rcode}

\textbf{Question:} What does power = 0.345 mean in plain English?

\vspace{1.5cm}

\subsection{Exercise 1.2: Using power.t.test()}

R has built-in functions for power analysis. You specify any 4 of the 5 parameters, and R solves for the missing one.

\begin{rcode}
# Calculate power for given sample size
result <- power.t.test(n = 25,           # Sample size
                       delta = 10,       # Effect size |mu1 - mu0|
                       sd = 30,          # Standard deviation
                       sig.level = 0.05, # Alpha
                       type = "one.sample",
                       alternative = "two.sided")
print(result)
cat("\nPower:", round(result$power, 3), "\n")

# Calculate required sample size for power = 0.80
result2 <- power.t.test(delta = 10,
                        sd = 30,
                        sig.level = 0.05,
                        power = 0.80,
                        type = "one.sample")
print(result2)
cat("\nRequired n:", ceiling(result2$n), "\n")

# Result: Need n = 143 subjects for 80% power
\end{rcode}

\textbf{Key insight:} Small samples have low power! Always do power analysis before collecting data.

\vspace{1cm}

\subsection{Exercise 1.3: Power Curves}

Visualize how power changes with sample size.

\begin{rcode}
# Create power curve
sample_sizes <- seq(10, 200, by=5)

# Calculate power for each sample size
powers <- sapply(sample_sizes, function(n) {
  power.t.test(n = n, delta = 10, sd = 30,
               sig.level = 0.05, type = "one.sample")$power
})

# Plot power curve
plot(sample_sizes, powers, type="l", lwd=2, col="blue",
     main="Power vs. Sample Size",
     xlab="Sample Size (n)",
     ylab="Power",
     ylim=c(0, 1))

# Add reference lines
abline(h=0.80, col="red", lty=2, lwd=1.5)
abline(v=143, col="red", lty=2, lwd=1.5)
text(143, 0.5, "n = 143", pos=4, col="red")
text(100, 0.83, "Power = 0.80", pos=4, col="red")
grid()

# Key observation: Power increases with n, but with diminishing returns
\end{rcode}

\textbf{Question:} Why does the power curve flatten out at large sample sizes?

\vspace{1.5cm}

\subsection{Exercise 1.4: Two-Sample Power}

For comparing two groups, use \texttt{type = "two.sample"}.

\begin{rcode}
# Sample size for two-sample t-test
# Want to detect 15 mmHg difference in blood pressure
# Assume SD = 20 mmHg in both groups
# Want 90% power

power.t.test(delta = 15,          # Difference to detect
             sd = 20,              # Pooled SD
             sig.level = 0.05,
             power = 0.90,
             type = "two.sample",  # Two independent groups
             alternative = "two.sided")

# Result: Need n = 34 per group (68 total)

# What if we only have n = 20 per group?
power.t.test(n = 20, delta = 15, sd = 20,
             sig.level = 0.05, type = "two.sample")
# Power = 0.724 (72.4%) - a bit low!
\end{rcode}

\section{Part 2: Two-Sample t-Tests (Independent Samples)}

Compare means of two independent groups.

\subsection{Exercise 2.1: Pooled Two-Sample t-Test}

\textbf{Study:} Compare tumor growth (mm) in control vs. treatment group.

\begin{rcode}
# Enter data
control <- c(7, 10, 9, 8, 7, 6, 8, 9, 12, 13)
treatment <- c(4, 6, 10, 8, 5, 3, 10, 8, 8, 10)

# Summary statistics
cat("Control group:\n")
cat("  Mean:", mean(control), "\n")
cat("  SD:", sd(control), "\n")
cat("  n:", length(control), "\n\n")

cat("Treatment group:\n")
cat("  Mean:", mean(treatment), "\n")
cat("  SD:", sd(treatment), "\n")
cat("  n:", length(treatment), "\n\n")

# Visualize data
boxplot(control, treatment,
        names=c("Control", "Treatment"),
        ylab="Tumor Growth (mm)",
        main="Tumor Growth: Control vs. Treatment",
        col=c("lightcoral", "lightblue"))

# Add points
stripchart(list(Control=control, Treatment=treatment),
           vertical=TRUE, method="jitter", add=TRUE,
           pch=19, col="darkgray")
\end{rcode}

\textbf{Observation:} Treatment group appears to have slightly lower growth, but there's overlap.

\newpage

\begin{rcode}
# Pooled two-sample t-test
# Assumes equal variances
t.test(control, treatment,
       var.equal = TRUE,         # Use pooled variance
       alternative = "two.sided")

# Output interpretation:
# t = 1.571
# df = 18 (n1 + n2 - 2 = 10 + 10 - 2)
# p-value = 0.134
# 95% CI for difference: (-0.56, 3.96) mm

# Conclusion: p = 0.134 > 0.05
# Not enough evidence to conclude treatment reduces tumor growth
\end{rcode}

\textbf{Interpretation:} Write a complete conclusion in context:

\vspace{2cm}

\subsection{Exercise 2.2: Welch's t-Test (Unequal Variances)}

Welch's test does not assume equal variances. It's the default in R.

\begin{rcode}
# Welch's two-sample t-test (default)
t.test(control, treatment,
       var.equal = FALSE)  # Or just omit this line

# Slightly different results:
# t = 1.571
# df = 17.78 (adjusted for unequal variances)
# p-value = 0.134

# For this data, results are nearly identical because
# variances are similar (2.234 vs 2.573)

# Check variance equality with F-test
var.test(control, treatment)
# F = 0.754, p = 0.638
# p > 0.05 so variances are not significantly different
# Either pooled or Welch's test is appropriate
\end{rcode}

\textbf{Rule of thumb:} When in doubt, use Welch's test (safer, default in R).

\vspace{1cm}

\subsection{Exercise 2.3: One-Sided Tests}

Sometimes we have a directional hypothesis.

\begin{rcode}
# One-sided test: Is control > treatment?
# (Does treatment reduce growth?)
t.test(control, treatment,
       alternative = "greater")  # Test if control > treatment

# p-value = 0.0671
# One-sided p-value is half of two-sided when t > 0

# Still not significant at alpha = 0.05,
# but closer (p = 0.067 vs 0.134)
\end{rcode}

\textbf{Note:} Only use one-sided tests when you have strong prior reason to expect direction!

\vspace{1cm}

\subsection{Exercise 2.4: Confidence Intervals}

The confidence interval tells us the plausible range for the true difference.

\begin{rcode}
# Extract confidence interval
result <- t.test(control, treatment, var.equal=TRUE)

# Access components
cat("Difference in means:", result$estimate[1] - result$estimate[2], "\n")
cat("95% CI:", result$conf.int[1], "to", result$conf.int[2], "\n")

# Since CI includes 0, not significant at alpha = 0.05
# If CI were (0.5, 4.0), would be significant

# For one-sided test, get one-sided CI
result_one <- t.test(control, treatment,
                     var.equal=TRUE,
                     alternative="greater")
cat("One-sided 95% CI:", result_one$conf.int[1], "to Inf\n")
\end{rcode}

\section{Part 3: Paired t-Tests}

Use paired tests when observations are matched or dependent.

\subsection{Exercise 3.1: Before-After Study}

\textbf{Study:} Blood pressure before and after medication in 10 patients.

\begin{rcode}
# Enter data
before <- c(142, 138, 150, 148, 135, 160, 155, 145, 152, 158)
after <- c(138, 132, 148, 140, 135, 152, 148, 142, 146, 150)

# Calculate differences
differences <- before - after

cat("Summary of differences:\n")
cat("  Mean:", mean(differences), "mmHg\n")
cat("  SD:", sd(differences), "\n")
cat("  n:", length(differences), "\n")

# Visualize
par(mfrow=c(1,2))

# Before-after plot
plot(before, after, pch=19, col="blue",
     xlab="Before (mmHg)", ylab="After (mmHg)",
     main="Blood Pressure: Before vs After")
abline(0, 1, col="red", lty=2)  # Line of no change
for(i in 1:10) {
  segments(before[i], before[i], before[i], after[i], col="gray")
}

# Histogram of differences
hist(differences, breaks=5, col="lightblue",
     main="Distribution of Differences",
     xlab="Change in BP (mmHg)")
abline(v=mean(differences), col="red", lwd=2)

par(mfrow=c(1,1))
\end{rcode}

\newpage

\begin{rcode}
# Paired t-test (Method 1: test differences)
t.test(differences, mu=0, alternative="greater")

# Output:
# t = 5.916, df = 9
# p-value = 0.0001236
# mean of differences = 5.2 mmHg

# Paired t-test (Method 2: paired argument)
t.test(before, after, paired=TRUE, alternative="greater")

# Conclusion: Strong evidence (p < 0.001) that medication
# reduces blood pressure by an average of 5.2 mmHg
\end{rcode}

\textbf{Interpretation:} Write a conclusion in context:

\vspace{2cm}

\subsection{Exercise 3.2: Why Pairing Matters}

Compare paired vs. independent analysis on the same data.

\begin{rcode}
# WRONG: Independent two-sample t-test
independent_result <- t.test(before, after, paired=FALSE)
cat("Independent t-test:\n")
cat("  t =", independent_result$statistic, "\n")
cat("  p-value =", independent_result$p.value, "\n\n")

# CORRECT: Paired t-test
paired_result <- t.test(before, after, paired=TRUE)
cat("Paired t-test:\n")
cat("  t =", paired_result$statistic, "\n")
cat("  p-value =", paired_result$p.value, "\n\n")

# Dramatic difference!
# Independent: p = 0.115 (not significant)
# Paired: p = 0.0002 (highly significant)

# Why? Compare standard errors
se_independent <- sqrt(var(before)/10 + var(after)/10)
se_paired <- sd(differences)/sqrt(10)

cat("SE (independent):", round(se_independent, 2), "\n")
cat("SE (paired):", round(se_paired, 2), "\n")
cat("Ratio:", round(se_independent/se_paired, 2), "\n")

# SE for paired is nearly 4x smaller!
# Pairing removes between-subject variability
\end{rcode}

\textbf{Lesson:} Always use paired test for paired data. Much more powerful!

\vspace{1cm}

\section{Part 4: One-Way ANOVA}

Compare means across 3 or more groups.

\subsection{Exercise 4.1: Drug Dose Study}

\textbf{Study:} Compare cholesterol reduction (mg/dL) for 3 drug doses.

\begin{rcode}
# Enter data
low <- c(10, 12, 8, 11, 9)
medium <- c(15, 18, 16, 20, 17)
high <- c(25, 28, 30, 26, 27)

# Create data frame (long format)
cholesterol <- c(low, medium, high)
dose <- factor(rep(c("Low", "Medium", "High"), each=5),
               levels=c("Low", "Medium", "High"))
data <- data.frame(cholesterol, dose)

# Summary statistics by group
aggregate(cholesterol ~ dose, data=data, FUN=mean)
aggregate(cholesterol ~ dose, data=data, FUN=sd)

# Visualize
boxplot(cholesterol ~ dose, data=data,
        col=c("lightblue", "lightgreen", "lightcoral"),
        main="Cholesterol Reduction by Dose",
        xlab="Dose", ylab="Reduction (mg/dL)")

# Add points
stripchart(cholesterol ~ dose, data=data,
           vertical=TRUE, method="jitter", add=TRUE,
           pch=19, col="darkgray")
\end{rcode}

\newpage

\subsection{Exercise 4.2: Check ANOVA Assumptions}

Before running ANOVA, verify assumptions.

\begin{rcode}
# 1. Normality: Check within each group
par(mfrow=c(1,3))
hist(low, main="Low Dose", xlab="Reduction", col="lightblue")
hist(medium, main="Medium Dose", xlab="Reduction", col="lightgreen")
hist(high, main="High Dose", xlab="Reduction", col="lightcoral")
par(mfrow=c(1,1))

# With small samples, hard to assess normality
# Q-Q plots are better
par(mfrow=c(1,3))
qqnorm(low, main="Low Dose"); qqline(low)
qqnorm(medium, main="Medium Dose"); qqline(medium)
qqnorm(high, main="High Dose"); qqline(high)
par(mfrow=c(1,1))

# 2. Equal variances (homogeneity)
# Bartlett's test (sensitive to non-normality)
bartlett.test(cholesterol ~ dose, data=data)

# Levene's test (more robust)
# install.packages("car")
library(car)
leveneTest(cholesterol ~ dose, data=data)

# p > 0.05 means equal variances (good!)

# Visual check: Boxplots should have similar spreads
# Our boxplots look okay
\end{rcode}

\textbf{Assumption check:} Are the ANOVA assumptions reasonably met for this data?

\vspace{2cm}

\subsection{Exercise 4.3: Perform ANOVA}

\begin{rcode}
# One-way ANOVA
model <- aov(cholesterol ~ dose, data=data)
summary(model)

# Output:
#             Df Sum Sq Mean Sq F value   Pr(>F)
# dose         2 985.73  492.87   149.4 7.77e-09 ***
# Residuals   12  39.60    3.30
# ---
# Signif. codes:  0 '***' 0.001 '**' 0.01 '*' 0.05 '.' 0.1 ' ' 1

# Interpretation:
# F(2, 12) = 149.4
# p < 0.001 (highly significant)
# Strong evidence that at least one dose differs

# Effect size: How much of variance explained by dose?
# R-squared (eta-squared)
ss_between <- 985.73
ss_total <- 985.73 + 39.60
r_squared <- ss_between / ss_total
cat("R-squared:", round(r_squared, 3), "\n")
# 96% of variance explained by dose!
\end{rcode}

\textbf{ANOVA conclusion:} There is strong evidence that mean cholesterol reduction differs across dose levels (F(2,12) = 149.4, p $<$ 0.001).

\vspace{1cm}

\subsection{Exercise 4.4: Post-Hoc Tests}

ANOVA tells us \textit{some} groups differ, but not \textit{which} ones. Use post-hoc tests.

\begin{rcode}
# Tukey's Honestly Significant Difference
TukeyHSD(model, conf.level=0.95)

# Output:
#   Tukey multiple comparisons of means
#     95% family-wise confidence level
#
# $dose
#                    diff       lwr      upr     p adj
# Medium-Low      7.2000   4.413   10.0870  0.000032
# High-Low       17.2000  14.413   19.9870  0.000000
# High-Medium    10.0000   7.213   12.7870  0.000001

# All three pairwise comparisons are significant!
# High > Medium > Low

# Visualize Tukey results
plot(TukeyHSD(model), las=1,
     main="Tukey HSD 95% Confidence Intervals")
abline(v=0, lty=2, col="red")

# If CI doesn't include 0, difference is significant
\end{rcode}

\textbf{Conclusion:} All three doses produce significantly different reductions. High dose is most effective, followed by medium, then low.

\vspace{1cm}

\subsection{Exercise 4.5: Alternative Post-Hoc Methods}

\begin{rcode}
# Bonferroni correction: Adjust p-values
pairwise.t.test(data$cholesterol, data$dose,
                p.adjust.method = "bonferroni")

# More conservative than Tukey for many comparisons
# For 3 groups (3 comparisons), similar results

# Holm's method (less conservative than Bonferroni)
pairwise.t.test(data$cholesterol, data$dose,
                p.adjust.method = "holm")

# No adjustment (for comparison - DON'T DO THIS!)
pairwise.t.test(data$cholesterol, data$dose,
                p.adjust.method = "none")
# Shows why we need correction!
\end{rcode}

\section{Part 5: Integrated Example}

Let's work through a complete analysis from start to finish.

\subsection{Exercise 5.1: Sleep Medication Study}

\textbf{Study:} Compare 4 sleep medications (A, B, C) vs. placebo. Measure hours of sleep.

\begin{rcode}
# Generate realistic data
set.seed(524)
placebo <- round(rnorm(8, mean=5.5, sd=1.2), 1)
med_a <- round(rnorm(8, mean=6.8, sd=1.1), 1)
med_b <- round(rnorm(8, mean=7.2, sd=1.0), 1)
med_c <- round(rnorm(8, mean=6.5, sd=1.3), 1)

# Create data frame
sleep_hours <- c(placebo, med_a, med_b, med_c)
medication <- factor(rep(c("Placebo", "Med_A", "Med_B", "Med_C"),
                         each=8))
sleep_data <- data.frame(sleep_hours, medication)

# View data
head(sleep_data, 10)
\end{rcode}

\newpage

\textbf{Step 1: Exploratory Analysis}

\begin{rcode}
# Summary statistics
aggregate(sleep_hours ~ medication, data=sleep_data, FUN=mean)
aggregate(sleep_hours ~ medication, data=sleep_data, FUN=sd)
aggregate(sleep_hours ~ medication, data=sleep_data, FUN=length)

# Visualizations
par(mfrow=c(1,2))

# Boxplots
boxplot(sleep_hours ~ medication, data=sleep_data,
        col=rainbow(4),
        main="Sleep Hours by Medication",
        xlab="Medication", ylab="Hours")

# Dotplots with means
stripchart(sleep_hours ~ medication, data=sleep_data,
           vertical=TRUE, pch=19, method="jitter",
           main="Individual Observations",
           xlab="Medication", ylab="Hours")
means <- tapply(sleep_data$sleep_hours,
                sleep_data$medication, mean)
points(1:4, means, pch=18, col="red", cex=2)
legend("topleft", legend="Group Mean",
       pch=18, col="red", pt.cex=2)

par(mfrow=c(1,1))
\end{rcode}

\textbf{Initial observations:} Write what you see from the plots.

\vspace{2cm}

\textbf{Step 2: Check Assumptions}

\begin{rcode}
# Normality
par(mfrow=c(2,2))
for(med in levels(sleep_data$medication)) {
  subset_data <- sleep_data$sleep_hours[sleep_data$medication == med]
  qqnorm(subset_data, main=med)
  qqline(subset_data)
}
par(mfrow=c(1,1))

# Equal variances
leveneTest(sleep_hours ~ medication, data=sleep_data)

# If p > 0.05, assumptions met
\end{rcode}

\textbf{Step 3: Perform ANOVA}

\begin{rcode}
# Fit ANOVA model
sleep_model <- aov(sleep_hours ~ medication, data=sleep_data)
summary(sleep_model)

# Extract key values
anova_table <- summary(sleep_model)[[1]]
f_stat <- anova_table$`F value`[1]
p_value <- anova_table$`Pr(>F)`[1]
df1 <- anova_table$Df[1]
df2 <- anova_table$Df[2]

cat("\nANOVA Results:\n")
cat("F(", df1, ",", df2, ") = ", round(f_stat, 2), "\n", sep="")
cat("p-value = ", round(p_value, 4), "\n", sep="")
\end{rcode}

\textbf{Step 4: Post-Hoc Tests (if significant)}

\begin{rcode}
# If ANOVA is significant (p < 0.05), do post-hoc
if(p_value < 0.05) {
  cat("\nANOVA is significant. Performing post-hoc tests...\n\n")

  # Tukey HSD
  tukey_results <- TukeyHSD(sleep_model)
  print(tukey_results)

  # Plot
  plot(tukey_results, las=1)
  abline(v=0, lty=2, col="red")
} else {
  cat("\nANOVA is not significant. No post-hoc tests needed.\n")
}
\end{rcode}

\textbf{Step 5: Write Conclusions}

Based on your analysis, write a complete conclusion answering:
\begin{enumerate}
\item Is there evidence that medications differ from placebo?
\item Which medication(s) are significantly better than placebo?
\item Which medication appears most effective?
\end{enumerate}

\vspace{3cm}

\section{Practice Problems}

\subsection{Problem 1: Power Analysis}

A researcher wants to test if a new teaching method improves test scores. Based on pilot data, the standard deviation is 12 points. The researcher wants to detect an improvement of 5 points.

\begin{enumerate}[label=\alph*)]
\item Calculate the required sample size for 80\% power at $\alpha = 0.05$ (two-sided).
\item What sample size would be needed for 90\% power?
\item If only 50 students are available, what power would the study have?
\end{enumerate}

\vspace{2cm}

\subsection{Problem 2: Two-Sample t-Test}

Compare weight loss (kg) between two diet programs:

Diet A: 3.2, 4.5, 2.8, 3.9, 4.1, 3.5, 4.3, 3.7

Diet B: 5.1, 4.8, 5.5, 4.9, 5.3, 5.0, 4.7, 5.2

\begin{enumerate}[label=\alph*)]
\item Visualize the data with boxplots.
\item Test if variances are equal using var.test().
\item Perform an appropriate two-sample t-test.
\item Calculate and interpret a 95\% CI for the difference.
\item Write a conclusion in context.
\end{enumerate}

\vspace{2cm}

\subsection{Problem 3: Paired t-Test}

Pain scores (0-10 scale) before and after treatment:

Before: 7, 8, 6, 9, 7, 8, 7, 6, 8, 7

After: 5, 6, 4, 7, 5, 6, 5, 4, 6, 5

\begin{enumerate}[label=\alph*)]
\item Calculate the differences.
\item Create a histogram of the differences.
\item Perform a paired t-test.
\item Compare to an independent two-sample t-test (incorrect!).
\item Explain why pairing matters for this data.
\end{enumerate}

\vspace{2cm}

\subsection{Problem 4: ANOVA}

Compare exam scores across 4 study methods:

Method 1: 75, 78, 72, 80, 76

Method 2: 82, 85, 80, 88, 83

Method 3: 70, 73, 68, 75, 71

Method 4: 85, 88, 83, 90, 86

\begin{enumerate}[label=\alph*)]
\item Create a data frame in long format.
\item Check ANOVA assumptions (normality, equal variances).
\item Perform one-way ANOVA.
\item If significant, perform Tukey's HSD.
\item Which study methods differ significantly?
\end{enumerate}

\vspace{2cm}

\section{\textcolor{aiblue}{AI-Enhanced Learning Tips}}

\begin{tcolorbox}[colback=aiblue!10,colframe=aiblue,title=Using AI Tools for Hypothesis Testing]

\textbf{Good prompts for AI assistance:}
\begin{itemize}
\item ``Calculate required sample size to detect 10-point difference with power 0.80 and SD=15''
\item ``Should I use paired or independent t-test for before/after measurements on same subjects?''
\item ``Generate R code for one-way ANOVA with 4 groups and Tukey post-hoc tests''
\item ``Explain the difference between Type I and Type II errors with a medical example''
\item ``How do I check homogeneity of variance assumption for ANOVA in R?''
\end{itemize}

\textbf{AI can help with:}
\begin{itemize}
\item Choosing the appropriate test
\item Power and sample size calculations
\item Checking assumptions
\item Interpreting output
\item Creating visualizations
\item Understanding when to use post-hoc tests
\end{itemize}

\textbf{Remember:}
\begin{itemize}
\item Always verify test assumptions before running
\item Don't confuse paired and independent samples
\item Check AI's test selection (paired vs. independent, one-sided vs. two-sided)
\item Power analysis should be done BEFORE data collection
\item Post-hoc tests only after significant ANOVA
\end{itemize}
\end{tcolorbox}

\section{Key Takeaways}

\begin{enumerate}
\item \textbf{Power:} Always calculate before study. Need $\geq$ 0.80. Increases with larger n, larger effect, smaller $\sigma$.
\item \textbf{Two-sample t-test:} Compare independent groups. Use pooled (equal var) or Welch's (unequal var).
\item \textbf{Paired t-test:} Much more powerful for dependent data. Always use for before/after or matched pairs.
\item \textbf{ANOVA:} Compare 3+ groups simultaneously. Controls Type I error better than multiple t-tests.
\item \textbf{Post-hoc:} Use Tukey's HSD after significant ANOVA to find which groups differ.
\item \textbf{Assumptions:} Independence, normality, equal variances. Check before testing!
\item \textbf{R Functions:}
  \begin{itemize}
  \item \texttt{power.t.test()}: Power and sample size
  \item \texttt{t.test()}: One-sample, two-sample, paired t-tests
  \item \texttt{aov()}: ANOVA
  \item \texttt{TukeyHSD()}: Post-hoc tests
  \item \texttt{var.test()}, \texttt{leveneTest()}: Check equal variances
  \end{itemize}
\end{enumerate}

\end{document}
